\documentclass[11pt]{article}

\usepackage{amsmath}
\usepackage{amsfonts}
\usepackage{amsthm,thmtools}
\usepackage{tikz}
\usepackage{enumerate}


\newcommand{\batsym}{\begin{tikzpicture}[scale=0.05]
	\draw[fill,black]   (-0.25,1.48) .. controls (-0.1,1.5) and (0.1,1.5) .. (0.25,1.48) -- (0.35,1.92) .. controls (0.425,1.8) and (0.41,1.3) .. (0.45,1.2) .. controls (0.6,1.05) and (1.96,1.05) .. (1.98,2.08) -- (5.93,2.08) .. controls (4.2,1.45) and (4,0.3) .. (4.2,-0.28) .. controls (2.4,-0.09) and (0.4,-0.5) .. (0,-2.052) .. controls (-0.4,-0.5) and (-2.4,-0.09) .. (-4.2,-0.28) .. controls (-4,0.3) and (-4.2,1.45) .. (-5.93,2.08) -- (-1.98,2.08) .. controls (-1.96,1.05) and (-0.6,1.05) .. (-0.45,1.2) .. controls (-0.41,1.3) and (-0.425,1.8) .. (-0.35,1.92) -- (-0.25,1.48);
	\end{tikzpicture}}
\newcommand{\QEB}{\begin{flushright}\batsym\end{flushright}}



\def\bold{\textbf}
\theoremstyle{plain}
\newtheorem{problem}{Problem}
\theoremstyle{remark}
\newtheorem*{idea}{Idea}

\renewcommand*{\proofname}{Solution}


\renewcommand{\qedsymbol}{$\batsym$}

\newcommand {\ds}{\displaystyle\sum}
\newcommand {\nnR} {\mathbb{R}_{\geq 0}}

\newcommand {\Mtt}[4] {\begin{bmatrix}#1 & #2\\ #3 & #4\end{bmatrix}}
\newcommand {\Ctt}[2] {\Mtt {#1}{-#2}{#2}{#1}}
\newcommand {\mbR} {\mathbb R}
\newcommand {\mcC} {\mathcal C}
\newcommand {\mbZ} {\mathbb Z}
\newcommand {\mbN} {\mathbb N}
\newcommand {\mbQ} {\mathbb Q}
\newcommand {\RnQ} {\mathbb R \setminus \mathbb Q}
\newcommand {\mcI} {\mathcal I}

\begin{document}
	\title{Math 3A03 - Tutorial 5 Questions - Winter 2019}
	\author{Nikolay Hristov}
	\date{February 11/13, 2019}
	\maketitle
	
		\begin{problem}
		Which of the following sets are countable:
		
		\begin{enumerate}[(a)]
			\item $\mbR \setminus \mbQ$.
			
			\item $\left\{ \frac 1 n : n \in \mbN \right\}$.
			
			\item $A\times B = \left\{ (a,b): a\in A, b\in B \right\}$, where A and B are countable.
			
			\item $\mbQ \cup \left(\sqrt 2 \mbQ \right)$, where $\sqrt 2 \mbQ = \left\{\sqrt 2 q: q\in \mbQ \right\}$.
		\end{enumerate}
	\end{problem}

	\begin{proof}
		(a) is not countable, the remaining sets are countable.
	\end{proof}

	\begin{problem}
		Let $A,B$ be countable sets, and $\mathcal C=\{a+bi: a \in A,b\in B\}$, where $i^2=-1$, prove that $\mathcal C$ is countable.
	\end{problem}

	\begin{proof}
		Since $A,B$ are countable, then we have two functions ${f_A(x):A\rightarrow \mbN,}\newline{f_B(x):B\rightarrow \mbN}$, which are both bijective. Suppose we have an element $a+bi\in \mcC$, define $f(a+bi)=2^{f_A(a)}3^{f_B(b)}$, which is a function mapping $\mcC$ to $\mbN$. We'd like to show this function is injective and surjective. 
		
		First let's show injectivity, suppose that $a+bi,c+di\in\mcC$, and $f(a+bi)=f(c+di)\implies 2^{f_A(a)}3^{f_B(b)}=2^{f_A(c)}3^{f_B(d)} $. Since $\gcd(2,3)=1$ we have that $f_A(a)=f_A(c)$ and $f_B(b)=f_B(d)$, further the functions $f_A$ and $f_B$ are bijective so they are also injective, and so $a=c$ and $b=d$, which shows injectivity for $f$.
		
		Surjectivity will be a tad tricky, because $f$ \textit{isn't} surjective on $\mbN$, it is however surjective on a subset of $\mbN$ which we will define in a minute. Just recall that (as you will prove in your assignment), a subset of a countable set is also countable, so having a bijection to a subset of the natural numbers is enough. How do we define this subset? Recall a surjection roughly means that you map to the entirety of set, so we'd like to use $\mcI = f(\mcC) = \{f(a+bi):a\in A, b\in B\}$, i.e. the image of the set $\mcC$ under the function $f$. 
		
		The function is still injective. For surjectivity, suppose that we have a $z\in \mcI$, we need to find an $a,b$ such that $z=f(a+bi)$, but by definition each element of the set is $f(a+bi)$ for some $a\in A$ and $b\in B$. Hence $f$ is injective and surjective as a map $f:\mcC \rightarrow \mcI\subset \mbN$, so it is a bijective map to $\mcI$. Since $\mcI$ is a subset of a countable set, by a problem in assignment 3 we have that $\mcI$ is countable. We can now conclude that $\mcC$ is countable, to be a bit more complete we can say that since $\mcI$ is countable there exists a bijective function $g:\mcI \rightarrow \mbN$, notice that $g\circ f:\mcC \rightarrow \mbN$ is a composition of bijective functions, which is also bijective, and hence there is a bijective function from $\mcC$ to $\mbN$, and hence $\mcC$ is countable.
	\end{proof}

	\begin{problem}
		Show that every interior point is an accumulation point.
	\end{problem}

	\begin{proof}
		
		Suppose that $x$ is an interior point of some set $E\subseteq \mbR$, then there is a $c>0$ such that $(x-c,x+c)\subseteq E$. Now suppose that $d>0$ is an arbitrary positive number. We'd like to show that $(E\setminus \{x\})\cap (x-d,x+d)$ is nonempty. Notice however that since $(x-c,x+c)\subseteq E $ then \[(x-d,x+d)\cap (x-c,x+c)\subseteq (x-d,x+d)\cap E.\]
		
		Further if we let $a = \min (c,d)$ then $(x-d,x+d)\cap (x-c,x+c)=(x-a,x+a)$, removing the point ${x}$ we end up with the union of two intervals, $(x-a,x)$ and $(x,x+a)$, hence
		
		\[(x-a,x)\cup (x,x+a)\subseteq (x-d,x+d)\cap (E\setminus {x}).\]
		
		Since the intervals are nonempty then $(x-d,x+d)\cap (E\setminus {x})$ contains at least one point.
		
		
	\end{proof}

	\begin{problem}
		Write the open interval $(0,2)$ as a union of closed sets. Can it be expressed as an intersection of closed sets?
	\end{problem}
	\begin{proof}
		$(0,2) = \displaystyle\cup_{n=1}^\infty \left[\frac 1 n, 2-\frac{1}{n}\right]$. To justify this note that whenever $0<x<2$ there is a natural number n, sufficiently large, such that $x\geq \frac 1 n$ and $x\leq 2 - \frac 1 n$, and hence $x\in [\frac 1 n, 2-\frac{1}{n}]$ for some n. To show that the boundary points aren't in the set, note that for every $n\in \mbN$ we have that $0< \frac 1 n$, similarly $2> 2-\frac 1 n$.
		
		This cannot be expressed as an intersection of closed sets as any intersection of closed sets is also closed (see class notes).
	\end{proof}
	\begin{problem}
		
		\begin{enumerate} [(a)]
		\item Find the closure, interior points, accumulation points, and boundary points of the set $S=[0,\sqrt 5] \cap \mbQ$. Is this set open? Is it closed?
		
		\item Show that $O=(0,3)$ is open. 
		
		\item Show that $C=[0,3]$ is closed.
		
	\end{enumerate}
	\end{problem}

	\begin{proof}
		
		Next tutorial!
		
		\iffalse
		(a) The closure of the set is $[0, \sqrt 5]$. The interior points of the set are $\emptyset$, since for any point q in the set and any $c>0$ we have points in $[q-c,q+c]$ which are not in the set since between any two real numbers there is an irrational number. The accumulation points of the set are $[0,\sqrt 5]$, since by density of the real numbers given a $c>0$ and $x\in [0,\sqrt 5]$ we can find a rational number in $[x-c,x+c]\cap (S\setminus \{x\})$. By similar reasoning $[0,\sqrt 5]$ are the boundary points of S. The set is not open since there are points in the set which are not interior points. S is not closed since it does not contain all of its accumulation points.
		
		(b) To show that the set is open we need to verify that each $x\in (0,3)$ is an interior point. Suppose that $x\in (0,3)$, then $0<x<3$. To verify this is an interior point we need to find a $c>0$ such that $(x-c,x+c)\subseteq (0,3)$. Note that if $0<c\leq\frac x 2$ then $x-c\geq x - \frac x 2 = \frac x 2 >0$, further if $0<c\leq \frac {3-x} 2$ then $x+c\leq \frac {3+x} 2 < 3$. Then if we pick $c=\min (\frac x 2, \frac {3-x}2)$ the requirements that $0<c\leq \frac x 2$ and $0<c\leq \frac {3-x} 2$ are satisfied, hence $0<x-c<x+c<3$, and $(x-c,x+c)\subseteq (0,3)$.
		
		(c) We have 2 options, one is to show that the complement $[0,3]^C = (-\infty,0)\cup (3,\infty)$ is open, or we can show that the set contains all of its accumulation points. Let's do the latter. First note that if $x\in [0,3]$, then every neighbourhood of x contains infinitely many points of the interval, so $x$ is an accumulation point. We need to show that there are no accumulation points outside of the set. i.e. we need to show that if $x>3$ or $x<0$ then $x$ is \textbf{not} an accumulation point of the set. What we need to prove is for such an $x$ there exists a $c>0$ such that $(x-c,x+c)\cap ([0,3]\setminus \{x\})=\emptyset$. Suppose that $x<0$ and let $c=\frac{|x|}{2}$, then $x+c<0$ as well, hence $(x-c,x+c)\cap ([0,3]\setminus \{x\})=\emptyset$ since the sets are disjoint (note that $\sup (x-c,x+c)=-\frac{|x|}{2}< 0=\inf [0,3]$). A similar argument shows that each $x>3$ is also not an accumulation point.
		
		\fi
	\end{proof}


	
\end{document}